\documentclass[11pt]{article}
\newcommand{\MODULETAG}{Module 2\,\textbar\,Sampling \& Data}
% ===== Toro Resource Estimation Course - shared PDF preamble =====
\usepackage[a4paper,margin=2.4cm,headsep=14pt]{geometry}
\usepackage{amsmath,amssymb}
\usepackage{lmodern}
\usepackage[T1]{fontenc}
\usepackage[utf8]{inputenc}
\usepackage{microtype}
\usepackage{graphicx}
\usepackage{booktabs}
\usepackage{array}
\usepackage{enumitem}
\usepackage{xcolor}
\usepackage{titlesec}
\usepackage{fancyhdr}
\usepackage[most]{tcolorbox}
\usepackage{listings}
\usepackage{caption}
\usepackage[hidelinks]{hyperref}

% ---- palette ----
\definecolor{navy}{HTML}{1F2A44}
\definecolor{navy2}{HTML}{2B3A5C}
\definecolor{bronze}{HTML}{C8702D}
\definecolor{bronzed}{HTML}{A85A1F}
\definecolor{teal}{HTML}{2A7D8C}
\definecolor{tealw}{HTML}{ECF4F5}
\definecolor{ink}{HTML}{26303F}
\definecolor{paper2}{HTML}{F3EDE1}
\definecolor{line}{HTML}{D8D4CB}
\definecolor{codebg}{HTML}{1C2433}
\definecolor{good}{HTML}{3F7D5A}
\definecolor{warn}{HTML}{B4451F}

\color{ink}
\setlength{\parindent}{0pt}
\setlength{\parskip}{6pt}
\linespread{1.06}

% ---- section styling ----
\titleformat{\section}{\Large\bfseries\color{navy}}{\thesection}{0.6em}{}
\titleformat{\subsection}{\large\bfseries\color{navy2}}{\thesubsection}{0.6em}{}
\titleformat{\subsubsection}{\normalsize\bfseries\color{bronzed}}{}{0em}{}
\titlespacing{\section}{0pt}{16pt}{6pt}
\titlespacing{\subsection}{0pt}{12pt}{4pt}

% ---- header / footer ----
\pagestyle{fancy}
\fancyhf{}
\renewcommand{\headrulewidth}{0.4pt}
\renewcommand{\footrulewidth}{0pt}
\fancyhead[L]{\footnotesize\color{navy2}\MODULETAG}
\fancyhead[R]{\footnotesize\color{navy2}Toro Tin Project\,\textbar\,ANRML26-137}
\fancyfoot[C]{\footnotesize\color{navy2}\thepage}

% ---- callout boxes ----
\newtcolorbox{keybox}[1][Key concept]{breakable,enhanced,
  colback=paper2,colframe=bronze,boxrule=0pt,leftrule=4pt,
  arc=2pt,left=10pt,right=10pt,top=8pt,bottom=8pt,
  fonttitle=\bfseries\sffamily\small\color{bronzed},title={#1}}
\newtcolorbox{torobox}[1][Toro application]{breakable,enhanced,
  colback=navy!6,colframe=navy,boxrule=0pt,leftrule=4pt,
  arc=2pt,left=10pt,right=10pt,top=8pt,bottom=8pt,
  fonttitle=\bfseries\sffamily\small\color{navy},title={#1}}
\newtcolorbox{pitbox}[1][Common pitfall]{breakable,enhanced,
  colback=warn!7,colframe=warn,boxrule=0pt,leftrule=4pt,
  arc=2pt,left=10pt,right=10pt,top=8pt,bottom=8pt,
  fonttitle=\bfseries\sffamily\small\color{warn},title={#1}}
\newtcolorbox{notebox}[1][Note]{breakable,enhanced,
  colback=tealw,colframe=teal,boxrule=0pt,leftrule=4pt,
  arc=2pt,left=10pt,right=10pt,top=8pt,bottom=8pt,
  fonttitle=\bfseries\sffamily\small\color{teal},title={#1}}
\newtcolorbox{defbox}[1][Definition]{breakable,enhanced,
  colback=good!7,colframe=good,boxrule=0pt,leftrule=4pt,
  arc=2pt,left=10pt,right=10pt,top=8pt,bottom=8pt,
  fonttitle=\bfseries\sffamily\small\color{good},title={#1}}
\newtcolorbox{objbox}{enhanced,colback=navy,colframe=navy,
  arc=3pt,left=12pt,right=12pt,top=10pt,bottom=10pt,
  coltext=white}

% ---- python code ----
\lstdefinestyle{py}{language=Python,
  basicstyle=\footnotesize\ttfamily\color{white},
  keywordstyle=\color{bronze},commentstyle=\color{line},
  stringstyle=\color{teal!60!white},backgroundcolor=\color{codebg},
  numberstyle=\tiny\color{line},showstringspaces=false,
  breaklines=true,frame=single,framerule=0pt,
  rulecolor=\color{codebg},xleftmargin=8pt,xrightmargin=4pt,
  aboveskip=10pt,belowskip=6pt,framexleftmargin=8pt,
  framexrightmargin=8pt,framextopmargin=6pt,framexbottommargin=6pt}
\lstset{style=py}

% ---- equation box ----
\newtcolorbox{eqbox}{enhanced,colback=white,colframe=line,
  boxrule=0.6pt,arc=2pt,left=8pt,right=8pt,top=4pt,bottom=4pt,
  before skip=10pt,after skip=10pt}

% ---- captions ----
\captionsetup{font={small,it},labelfont={bf,color=navy},labelsep=endash}

% ---- module title macro ----
\newcommand{\moduletitle}[3]{%
  {\sffamily\bfseries\footnotesize\color{bronze}\MakeUppercase{#1}\par}
  \vspace{2pt}
  {\sffamily\bfseries\LARGE\color{navy}#2\par}
  \vspace{3pt}
  {\small\itshape\color{navy2}#3\par}
  \vspace{4pt}{\color{line}\hrule height 1.5pt}\vspace{12pt}}

\begin{document}

\moduletitle{Module 2 of 6}{Sampling, Data and Database Integrity}{An
estimate can never be better than the samples it is built from}

This module covers the part of the value chain that quietly determines
the accuracy of everything downstream --- and the part most often done
carelessly. We derive the theory of sampling error, then apply it to the
real integrity problems sitting inside the Toro database now.

\begin{objbox}
{\sffamily\bfseries\color{bronze}LEARNING OBJECTIVES}
\begin{itemize}[leftmargin=14pt,itemsep=2pt]
\item Quantify the fundamental sampling error using Gy's theory and
explain why particle size dominates it.
\item Design appropriate sampling for an alluvial placer and for a
primary hard-rock target.
\item Structure a drillhole / pit database into its four core tables.
\item Determine bulk density correctly and quantify its effect on
declared tonnage.
\item Build a QAQC system from standards, blanks and duplicates.
\item Audit the Toro dataset and composite it correctly.
\end{itemize}
\end{objbox}

\section{The error budget}

A Toro pit bulk sample is roughly 100--200 kg; a future mining operation
will process several million tonnes. Every tonne mined will have its
grade \emph{predicted} from a sampled fraction of the order of one part
in ten million. The total error in a final block estimate is a sum of
variances from independent stages:

\begin{eqbox}
\[ \sigma^2_{\text{total}} = \sigma^2_{\text{sampling}}
   + \sigma^2_{\text{sub-sampling}} + \sigma^2_{\text{assay}}
   + \sigma^2_{\text{estimation}} \tag{2.1}\]
\end{eqbox}

\begin{keybox}[The order that matters]
Modules 4 and 5 reduce only $\sigma^2_{\text{estimation}}$. If the
sampling and assay variances are large, no kriging plan can rescue the
estimate --- you would be interpolating noise with great precision.
Sampling quality is logically prior to geostatistics.
\end{keybox}

\section{Gy's theory of sampling}

Pierre Gy gave the industry a quantitative theory of sampling error. Its
centrepiece is the \textbf{Fundamental Sampling Error} (FSE): the
irreducible error that remains even with a perfect protocol, arising
because ore is made of discrete, heterogeneous particles.
\emph{Constitution heterogeneity} is intrinsic and cannot be mixed away;
\emph{distribution heterogeneity} is the spatial arrangement and can be
reduced by homogenisation. The FSE flows from constitution heterogeneity.
For a low-grade ore sampled correctly:

\begin{eqbox}
\[ \sigma^2_{\text{FSE}} = \left(\frac{1}{M_S}-\frac{1}{M_L}\right)
   f\,g\,c\,\ell\,d^{\,3} \tag{2.2}\]
\end{eqbox}

Here $M_S$ is the sample mass, $M_L$ the lot mass, $d$ the nominal top
particle size (note the \emph{cube}), $f$ the shape factor ($\approx
0.5$), $g$ the granulometric factor ($\approx 0.25$), $c$ the
mineralogical factor and $\ell$ the liberation factor. The mineralogical
factor carries the grade dependence:

\begin{eqbox}
\[ c = \frac{1-a_L}{a_L}\big[(1-a_L)\rho_M + a_L\rho_g\big]
   \;\xrightarrow[a_L\to 0]{}\; \frac{\rho_M}{a_L} \tag{2.3}\]
\end{eqbox}

Two consequences must shape how you sample. \textbf{Particle size
dominates, because of the cube}: halving the top size by crushing cuts
the sampling variance to one eighth --- which is why a sampling protocol
is a staged sequence of crush, split, crush, split. \textbf{Low grade
plus a dense, coarse mineral is the worst case}: since $c \propto 1/a_L$,
error explodes as grade falls. Cassiterite is dense
($\rho\approx6.9$~g/cm$^3$), often coarse and frequently liberated --- it
behaves like a coarse-gold ore, where a few heavy grains carry a
disproportionate share of the metal.

\begin{notebox}[Gy's theory and the nugget effect]
The FSE is a variance that exists even between two samples taken at the
\emph{same} place. On a variogram (Module 4) that same-location variance
appears as the intercept at zero distance --- the \textbf{nugget
effect}. A large part of a tin deposit's nugget is not geology; it is
Gy's fundamental sampling error plus assay error. Improving the sampling
protocol literally lowers the nugget you fight in Module 4.
\end{notebox}

Beyond the FSE, Gy catalogued \emph{avoidable} errors: the grouping and
segregation error (controlled by taking many small increments) and the
materialisation errors of delimitation, extraction and preparation. These
produce \emph{bias}, which is far more dangerous than random error
because no amount of repeat sampling reveals it.

\section{Sampling the two Toro orebodies}

The two Toro targets are sampled by different methods because they are
different estimation problems.

\subsection{The alluvial target: a volume-grade problem}

Alluvial cassiterite sits in wash gravels beneath barren overburden. The
economic quantity is a concentration of heavy mineral per unit volume,
reported in \textbf{grams of tin per cubic metre}. A pit is dug to the
wash layer, a measured bulk sample of gravel is washed and panned to a
heavy-mineral concentrate (HMC), the HMC is weighed and assayed. The
in-situ grade is a \emph{chain} of factors:

\begin{eqbox}
\[ g_{\,\mathrm{Sn}}\;\Big[\tfrac{\text{g}}{\text{m}^3}\Big]
   = \underbrace{\frac{M_{\text{HMC}}}{M_{\text{bulk}}}}_{\text{yield}}
   \times\; a_{\text{Sn}} \;\times\; \rho_{\text{gravel}}
   \;\times\; 10^{6} \tag{2.4}\]
\end{eqbox}

\begin{pitbox}[The most important point in this module]
The Toro XRF dataset reports the grade \emph{of the panned heavy-mineral
concentrate}, not the grade of the ground. A sample reading ``3.97\%
SnO$_2$'' describes a kilogram of concentrate, not a cubic metre of
gravel. To turn it into a resource grade you must multiply by the HMC
yield and the bulk density, as in equation (2.4). Treating concentrate
assays as in-situ grades would overstate the deposit by one to two orders
of magnitude --- the classic, fatal error in alluvial tin reporting.
\end{pitbox}

The chemistry confirms it: raw alluvium would assay 70--90\% SiO$_2$, but
the Toro field samples average roughly 32\% SiO$_2$ and 29\% Fe$_2$O$_3$
--- these are heavy-mineral concentrates.

\subsection{The primary target: a grade-tonnage problem}

The primary greisen and pegmatite target is sampled like any hard-rock
deposit, in \textbf{\% Sn} of rock in place: diamond core drilling (the
definitive method), reverse-circulation drilling, channel sampling of
veins, and pitting of weathered cover. Two cautions: \textbf{true width}
(grade applies to a thickness perpendicular to the structure, not the
down-hole interval) and \textbf{refractory chemistry} (cassiterite
resists aqua-regia digestion, so analyses must specify four-acid
digestion or lithium-borate fusion XRF --- which the ANRML26-137
laboratory correctly used).

\section{The drillhole and pit database}

Estimation software consumes a relational database of four core tables,
linked by a unique hole or pit identifier: \textbf{Collar} (location and
depth), \textbf{Survey} (the hole path), \textbf{Assay} (grade per
interval), and \textbf{Lithology} (the geological log that defines
domains). Every interval has a from/to depth; intervals never overlap or
gap; coordinates use one declared system (for Toro, WGS84 UTM Zone 32N).
The Toro \texttt{Priority Drilling Points} file is already a well-formed
collar table in embryo --- 13 collars at Bishiwai, Yadabongo and Kadade.

\section{Bulk density --- the silent half of every tonne}

Tonnage is volume times density, and a proportional error in density is a
proportional error in tonnage and contained metal, one-for-one:

\begin{eqbox}
\[ T = V \times \rho \qquad\Longrightarrow\qquad
   \frac{\Delta T}{T} = \frac{\Delta\rho}{\rho} \tag{2.5}\]
\end{eqbox}

Density must be measured by domain and weathering state, not assumed. For
the unconsolidated Toro wash gravel, core displacement is impossible: an
in-situ method (a measured excavation, weighed, with a surveyed void
volume) is required, recording the natural moisture state. For the
primary granite, fresh rock is near 2.6--2.7~t/m$^3$ but the weathered
cupola hosting near-surface tin can be far lower.

\section{QAQC: deciding whether the database can be trusted}

QAQC separates two properties of the data. \textbf{Accuracy} is freedom
from bias; \textbf{precision} is repeatability. They are independent, and
demonstrated by three control instruments: \textbf{Certified Reference
Materials} test accuracy, \textbf{blanks} test contamination, and
\textbf{duplicates} test precision --- each inserted blind at roughly one
in twenty. Precision is read from duplicate pairs as the relative
difference:

\begin{eqbox}
\[ \mathrm{RD} = \frac{|x_1-x_2|}{\tfrac{1}{2}(x_1+x_2)}\times 100\%
   \tag{2.6}\]
\end{eqbox}

\begin{pitbox}[What Toro's data cannot yet tell you]
The Toro file contains paired records at identical coordinates --- for
example DBP-033A and DBP-033A$'$, reading 7.59\% and 2.28\% SnO$_2$, a
relative difference above 100\%. But you cannot tell \emph{why}: the
suffixes might mean two stratigraphic layers (real geology) or a re-split
(poor precision). The reconnaissance sampling had no QAQC structure, so
the explanations are confounded. That is why a designed QAQC protocol ---
blind, randomly inserted duplicates --- must run through the next campaign
from the first sample.
\end{pitbox}

\section{Worked example: auditing the Toro database}

The file \texttt{Motimose\_GeoChemData.csv} holds 40 XRF records. Some
are not in-situ samples at all: identifiers carrying mineral-separation
tags (Nb, Fe, Zr, SnO2+SiO2) are laboratory heavy-mineral fractions. The
record \texttt{DBP-010 SnO2+SiO2-2} assays 65.4\% SnO$_2$ --- a lab-made
cassiterite concentrate that exists nowhere in the ground. It must be
filtered out.

\begin{lstlisting}
import pandas as pd, numpy as np
geo = pd.read_csv("Motimose_GeoChemData.csv")

# IDs carrying a mineral-separation tag are laboratory fractions,
# NOT in-situ field samples - they must not enter grade statistics.
sep_tags = ["Nb", "Fe", "Zr", "SnO2+SiO2", "FeB", "SnO2 (2)"]
def is_separation(sid):
    s = str(sid)
    return any(t in s for t in sep_tags) and len(s.split()) > 1

geo["record_type"] = np.where(
    geo["Lab_Sample_ID"].apply(is_separation),
    "lab_separation", "field_sample")
# field_sample 32 ; lab_separation 8
\end{lstlisting}

Eight of 40 records are laboratory fractions; 32 are genuine field
samples. Those 32 still over-represent pits logged with two layers, so we
\textbf{composite} to one grade per location:

\begin{lstlisting}
fs = geo[geo.record_type == "field_sample"].copy()
comp = (fs.groupby("Field_Pit_ID")
          .agg(SnO2=("SnO2", "mean"))   # length-weight if logged
          .reset_index())
# 32 field samples -> 26 pit locations
\end{lstlisting}

\begin{keybox}[The audit result]
The raw ``40-sample'' dataset resolves into \textbf{26 pit locations}
with a defensible in-situ concentrate grade. That number --- not 40 ---
is the real spatial dataset for Modules 3--5. Proper compositing should
be length-weighted; recording sampled thickness on every future pit is a
prerequisite for a correct composite.
\end{keybox}

\vspace{4pt}{\color{line}\hrule}\vspace{4pt}
{\small\itshape Module 2 of the Toro Resource Estimation course. Next:
Module 3, Exploratory Data Analysis and the Statistics of Grade.}

\end{document}
